\subsection{Causal Graphs, Interventions, and Identifiability}

Causal learning asks which predictions remain valid under interventions, not merely under passive observation.

\begin{definition} \label{def:sl-structural-causal-model} \hlt{Structural Causal Model}\\
A structural causal model specifies variables by assignments
\[
X_j=f_j(X_{\operatorname{pa}(j)},U_j),
\]
where the exogenous variables $U_j$ represent disturbances.  The associated directed graph records direct causal inputs.
\end{definition}

\begin{definition} \label{def:sl-intervention-do-operator} \hlt{Intervention}\\
An intervention $\operatorname{do}(X_j=x)$ replaces the structural assignment for $X_j$ by the constant $x$ and leaves the remaining assignments unchanged.  The resulting law is an interventional distribution.
\end{definition}

\begin{definition} \label{def:sl-causal-identifiability} \hlt{Identifiability}\\
A causal quantity is \hlt{identifiable} if it is uniquely determined by the observational distribution together with the assumed causal graph or model class.
\end{definition}

\begin{remark} \label{rem:sl-causal-versus-predictive} \hlt{Causal and Predictive Targets}\\
A predictor can be accurate under the training distribution while failing after an intervention.  Causal assumptions specify which conditional relationships are stable enough to transport across regimes.
\end{remark}
